\section{Output-dependent truncations and energy jumps}\label{sec:jumps}

We now allow each entry of one edge array to have its own truncation
interval.  Between endpoint scales, the heat estimate from
Section~\ref{sec:heat} still applies.  The new task is to bound the total
increase in energy when entries switch on or off.  We will in fact bound
the sum of the absolute energy jumps.  The estimate depends on the
original arrays alone, because every entry switches at most twice.
We first establish two inputs to this bound: estimates for Gaussian
averages that are uniform in the scale, and a formula for changes in
the matrix energy that remains valid when the rank changes.

\subsection{Maximal averages on a line}

For a locally integrable function on $\R$, and a sequence on $\mathbb Z$,
respectively, define the centered maximal functions
\[
 M_{\R}f(x)=\sup_{r>0}\frac1{2r}\int_{x-r}^{x+r}|f(t)|\,\dd t,
 \qquad
 M_{\mathbb Z}f(n)=\sup_{r\in\mathbb Z_{\ge0}}
       \frac1{2r+1}\sum_{|m-n|\le r}|f(m)|.
\]
Thus discrete averages include singletons.  We use Lebesgue measure on
$\R$ and counting measure on $\mathbb Z$.
The estimates below are versions of the classical Hardy--Littlewood
maximal inequalities~\cite[Section~IV]{HardyLittlewood1930}; we include
the centered covering argument and the Gaussian comparison with the
constants needed here.

\Needspace{3\baselineskip}
\begin{lemma}[One-dimensional maximal estimates]\label{lem:line-maximal}
On either of these measure spaces, write $M$ for the corresponding
centered maximal operator.  For $f\in L^1$ and $\rho>0$,
\[
 \mu\{Mf>\rho\}\le\frac3\rho\|f\|_1.
\]
For every $1<q<\infty$ and $f\in L^q$,
\begin{equation}\label{eq:line-maximal-strong}
 \|Mf\|_q^q\le
 \frac{6q\,2^{q-1}}{q-1}\|f\|_q^q.
\end{equation}
Moreover, there is an absolute constant $C_g$ such that, whenever
$h>0$, $\sigma\ge h^2$ and $n\in\mathbb Z$,
\begin{equation}\label{eq:discrete-gaussian-maximal}
 \sum_{m\in\mathbb Z}h g_\sigma((m-n)h)|f(m)|
 \le C_g M_{\mathbb Z}f(n).
\end{equation}
\end{lemma}

\begin{proof}
For a finite collection of intervals, repeatedly choose an interval of
largest radius and discard all intervals meeting it.  The chosen
intervals are disjoint.  Every discarded interval has radius no larger
than that of the interval which removed it and is contained in its
triple-radius enlargement.  This enlargement has three times the
original length on $\R$.  On $\mathbb Z$, its cardinality is
$6r+1\le3(2r+1)$, including $r=0$.

For each point in a finite subset of $\{M_{\mathbb Z}f>\rho\}$, choose
a witnessing interval whose average exceeds $\rho$.  The preceding
selection bounds the size of that subset by three times the sum of
the cardinalities of the chosen intervals, and hence by
$3\|f\|_1/\rho$.  Taking the supremum over finite subsets proves the
discrete weak estimate.  On $\R$, fixed-radius averages of an
$L^1$ function are continuous, so $\{M_{\R}f>\rho\}$ is open.
Apply the same selection to a finite witnessing cover of any compact
subset and then use inner regularity to obtain the continuous estimate.

For $f\in L^q$, split $f$ into its parts on $\{|f|\le\rho/2\}$ and
$\{|f|>\rho/2\}$.  The first part has maximal function at most
$\rho/2$, and the second is integrable.  The weak estimate gives
\[
 \mu\{Mf>\rho\}
 \le\frac6\rho\int_{|f|>\rho/2}|f|\,\dd\mu.
\]
The distribution formula and Tonelli's theorem now yield
\[
 \begin{aligned}
 \|Mf\|_q^q
 &\le6q\int |f(x)|
          \int_0^{2|f(x)|}\rho^{q-2}\,\dd\rho\,\dd\mu(x)\\
 &=\frac{6q\,2^{q-1}}{q-1}\|f\|_q^q.
 \end{aligned}
\]
These formulas are valid first with a bounded range of distribution
parameters and then by monotone convergence; finiteness of $Mf$ is
not assumed in advance.  Both maximal functions are measurable: the
discrete supremum is countable, and continuous averages are continuous
in their radius, so positive rational radii suffice on $\R$.

Finally put $a=\sqrt\sigma/h\ge1$.  In the sum in
\eqref{eq:discrete-gaussian-maximal}, the kernel is
\[
 h g_\sigma((m-n)h)
   =\frac1{\sqrt{2\pi}a}
       \exp\left(-\frac{(m-n)^2}{2a^2}\right).
\]
Split the integers into $|m-n|\le a$ and the annuli
$2^{j-1}a<|m-n|\le2^ja$, $j\ge1$.  The sum of $|f|$ over
$|m-n|\le2^ja$ is at most $3\cdot2^ja\,M_{\mathbb Z}f(n)$.
Bounding the Gaussian by its maximum on each region proves the claim
with, for example,
\[
 C_g=\frac3{\sqrt{2\pi}}
 \left(1+\sum_{j\ge1}2^j e^{-4^{j-1}/2}\right)<\infty.
\]
All constants are independent of the mesh and of the support of $f$.
\end{proof}

\subsection{Changing a matrix entry}

At a truncation endpoint the rank of a star can change.  We therefore
need a first-derivative formula that does not require fixed support.
For the sum of cubes of the singular values, this formula is a case
of the general gradient formula for unitarily invariant matrix
functions~\cite[Theorem~3.1]{Lewis1995}. We give a direct regularization
proof and then establish the entrywise bound used in the jump estimate.

\begin{lemma}[A derivative through rank changes]\label{lem:rank-change}
For finite rectangular complex matrices $R_-,R_+$ of the same size,
put $R(z)=(1-z)R_-+zR_+$ and $D=R_+-R_-$.  Then
\begin{equation}\label{eq:rank-change-derivative}
 \begin{aligned}
 &\Phi(R_+^*R_+)-\Phi(R_-^*R_-)\\
 &\quad=3\int_0^1\operatorname{Re}
   \langle R(z)\sqrt{R(z)^*R(z)},D\rangle_{\HS}\,\dd z.
 \end{aligned}
\end{equation}
For every rectangular matrix $R$ and every entry $(i,j)$,
\begin{equation}\label{eq:entry-gradient}
 |(R\sqrt{R^*R})_{ij}|
 \le\|\operatorname{row}_iR\|_2\,
       \|\operatorname{col}_jR\|_2.
\end{equation}
\end{lemma}

\begin{proof}
For $\eta>0$, differentiate
$\tr((R(z)^*R(z)+\eta I)^{3/2})$ using
Lemma~\ref{lem:hessian}.  Its derivative is
\[
 3\operatorname{Re}
 \langle R(z)\sqrt{R(z)^*R(z)+\eta I},D\rangle_{\HS}.
\]
Spectral calculus gives
\[
 \|\sqrt{R^*R+\eta I}-\sqrt{R^*R}\|_{\op}\le\sqrt\eta.
\]
Since $R(z)$ stays in a bounded set, the derivatives converge uniformly
in $z$, and the endpoint energies converge.  Integrating the
regularized identity before letting $\eta\downarrow0$ proves
\eqref{eq:rank-change-derivative}.

Cauchy--Schwarz in the matrix-product formula bounds the left side of
\eqref{eq:entry-gradient} by the norm of row $i$ of $R$ times the norm
of column $j$ of $\sqrt{R^*R}$.  The squared norm of that column is
$(R^*R)_{jj}$, which is also the squared norm of column $j$ of $R$.
This proves \eqref{eq:entry-gradient}, including zero columns and
singular matrices.
\end{proof}

\subsection{The total jump bound}

Return to the finite arrays and Gaussian weights of
Section~\ref{sec:heat}, and write $\dd\pi=\dd\pi_0$ on $\Pi_0$.
Independently for every pair $(i,j)$, choose
finite endpoints
\[
 h\le\varepsilon_{ij}<\rho_{ij},
 \qquad
 A_0^s(i,j)=A_0(i,j)
       \mathbf1_{\{\varepsilon_{ij}^2<s<\rho_{ij}^2\}}.
\]
The chosen endpoints are fixed independently of the Gaussian centers.
Keep $A_1,A_2$ fixed and form the weighted matrices and costs using
$A_0^s$ in place of $A_0$.  The arrays are complex; any phase at an
output pair may be included in the fixed value $A_0(i,j)$.
Set
\[
 s_- =\min_{i,j}\varepsilon_{ij}^2,
 \qquad s_+=\max_{i,j}\rho_{ij}^2.
\]

\begin{proposition}[Dissipation with output-dependent endpoints]
\label{prop:masked-dissipation}
For every $h>0$, $N\ge0$, complex arrays $A_v$ on
$\{-N,\ldots,N\}^2$, and endpoint choices as above, let
$\alpha$ be any permutation of $(1,11/10,11/10)$ and use variances
$\sigma_v=\alpha_vs$.  The corresponding masked costs satisfy
\begin{equation}\label{eq:masked-dissipation}
 \int_{s_-}^{s_+}\int_{\Pi_0}
       \sum_{v=0}^2 V_v(p,s)\,\dd\pi\,\dd s
 \le C\mathcal S(A).
\end{equation}
The constant is independent of the arrays, the mesh, the endpoints,
and the number of distinct endpoint values.
\end{proposition}

\begin{proof}
Write $J(s)=\int_{\Pi_0}\sum_ve_v(p,s)\,\dd\pi$ for the masked
energy, and put $\Delta J(\tau)=J(\tau+)-J(\tau-)$ at an endpoint.
We first estimate its jump at a single scale, then sum those
estimates and telescope the heat inequality.

\paragraph{The two stars at a jump.}
At an endpoint scale $\tau$, hold the weights at $s=\tau$ and let
\[
 \delta A(i,j)=A_0^{\tau+}(i,j)-A_0^{\tau-}(i,j).
\]
The superscripts denote one-sided masks.  Interpolate all changed
entries at once between these two arrays.  Each interpolated entry
has absolute value at most $|A_0(i,j)|$, even when some entries open
and others close at the same scale.

For fixed $i,j,p$, define nonnegative numbers using the original,
unmasked arrays:
\begin{equation}\label{eq:jump-slice-quantities}
 \begin{aligned}
 X^2&=\sum_m|A_0(i,m)|^2w_1(m),&
 Y^2&=\sum_m|A_0(m,j)|^2w_0(m),\\
 H^2&=\sum_m|A_2(m,i)|^2w_2(m),&
 K^2&=\sum_m|A_1(j,m)|^2w_2(m).
 \end{aligned}
\end{equation}
In $R_0=[W_0\;W_2^*]$, a changing entry has row norm at most
$\sqrt{w_0(i)}\sqrt{X^2+H^2}$, column norm at most
$\sqrt{w_1(j)}Y$, and derivative modulus
$\sqrt{w_0(i)w_1(j)}|\delta A(i,j)|$.
In $R_1=[W_1\;W_0^*]$, its adjoint entry has row norm at most
$\sqrt{w_1(j)}\sqrt{Y^2+K^2}$ and column norm at most
$\sqrt{w_0(i)}X$, with the same derivative modulus.
Star $R_2=[W_2\;W_1^*]$ is unchanged.  Lemma~\ref{lem:rank-change}
and $\sqrt{x^2+y^2}\le x+y$ therefore give
\begin{equation}\label{eq:jump-entry-bound}
 \begin{aligned}
 |\Delta e_0|+|\Delta e_1|
 \le3\sum_{i,j}|\delta A(i,j)|w_0(i)w_1(j)
       \bigl[(X+H)Y+(Y+K)X\bigr].
 \end{aligned}
\end{equation}
This estimate requires no assumption on the ranks during interpolation.

\paragraph{A probability measure on the center plane.}
For each fixed $i,j$, the measure
\[
 \dd\nu_{ij}=h^{-2}w_0(i)w_1(j)\,\dd\pi
\]
is a probability measure.  Indeed, in the free coordinates $(p_0,p_1)$,
its density is
$g_{\sigma_0}(u_i-p_0)g_{\sigma_1}(u_j-p_1)$.
Thus $p_0,p_1$ are independent Gaussian variables with means
$u_i,u_j$ and variances $\sigma_0,\sigma_1$, while
$p_2=-p_0-p_1$.  Adding Gaussian variances gives
\begin{equation}\label{eq:jump-gaussian-expectations}
 \begin{aligned}
 \mathbb E_{\nu_{ij}}w_0(m)
     &=h g_{2\sigma_0}(u_m-u_i),\\
 \mathbb E_{\nu_{ij}}w_1(m)
     &=h g_{2\sigma_1}(u_m-u_j),\\
 \mathbb E_{\nu_{ij}}w_2(m)
     &=h g_{\sigma_0+\sigma_1+\sigma_2}(u_m+u_i+u_j).
 \end{aligned}
\end{equation}
For the last identity, $p_0+p_1$ has mean $u_i+u_j$ and variance
$\sigma_0+\sigma_1$, and averaging $g_{\sigma_2}(u_m+p_0+p_1)$
adds the remaining variance.  Each identity also follows directly
by completing the square in its Gaussian integral.

Extend the arrays by zero to $\mathbb Z^2$ and define
\begin{equation}\label{eq:jump-starred-majorants}
 \begin{aligned}
 X_*^2(i,j)&=M_{\mathbb Z}(|A_0(i,\cdot)|^2)(j),\\
 Y_*^2(i,j)&=M_{\mathbb Z}(|A_0(\cdot,j)|^2)(i),\\
 H_*^2(i,j)&=M_{\mathbb Z}(|A_2(\cdot,i)|^2)(-i-j),\\
 K_*^2(i,j)&=M_{\mathbb Z}(|A_1(j,\cdot)|^2)(-i-j).
 \end{aligned}
\end{equation}
Every variance in \eqref{eq:jump-gaussian-expectations} is at least
$h^2$.  Lemma~\ref{lem:line-maximal} consequently bounds
$\mathbb E X^2,\mathbb E Y^2,\mathbb E H^2,\mathbb E K^2$ by
$C_g$ times the corresponding starred squares.  Cauchy--Schwarz
under the probability $\nu_{ij}$ gives, in particular,
\[
 \mathbb E(XY)\le C_gX_*Y_*,\qquad
 \mathbb E(HY)\le C_gH_*Y_*,\qquad
 \mathbb E(KX)\le C_gK_*X_*.
\]
No independence of these four random variables is needed.
The starred quantities depend only on the original arrays and the
output indices, not on the scale $\tau$.

\paragraph{Summing the jumps.}
Integrating \eqref{eq:jump-entry-bound} over the centers and using
$w_0(i)w_1(j)\,\dd\pi=h^2\,\dd\nu_{ij}$ yields
\[
 |\Delta J(\tau)|\le3C_g\sum_{i,j}h^2|\delta A(i,j)|
       (2X_*Y_*+H_*Y_*+K_*X_*).
\]
Every entry opens once and closes once.  Therefore its total absolute
change, summed over any subset of endpoint scales, is at most
$2|A_0(i,j)|$.  We obtain
\begin{equation}\label{eq:total-jump-majorant}
 \sum_\tau|\Delta J(\tau)|
 \le6C_g\sum_{i,j}h^2|A_0(i,j)|
       (2X_*Y_*+H_*Y_*+K_*X_*).
\end{equation}
The sum may include all endpoints or just the interior ones.

To bound this expression, set
\[
 n_v=\left(\sum_{i,j}h^2|A_v(i,j)|^3\right)^{1/3},
 \qquad \|f\|_{3,h}^3=\sum_{i,j\in\mathbb Z}h^2|f(i,j)|^3.
\]
The $q=3/2$ case of \eqref{eq:line-maximal-strong}, applied on slices,
gives
\begin{equation}\label{eq:jump-slice-norms}
 \|X_*\|_{3,h},\ \|Y_*\|_{3,h}\le Cn_0,
 \qquad \|H_*\|_{3,h}\le Cn_2,
 \qquad \|K_*\|_{3,h}\le Cn_1.
\end{equation}
For $H_*$, fix $i$ and change the summation index from $j$ to $-i-j$;
this is a bijection of $\mathbb Z$.  For $K_*$, fix $j$ and change
$i$ to $-i-j$.  Thus both estimates are ordinary slice estimates,
with no multiplicity or mesh factor.  Restricting the sums back to
the original finite square can only decrease them.

H\"older's inequality in \eqref{eq:total-jump-majorant}, followed by
\eqref{eq:jump-slice-norms}, proves
\begin{equation}\label{eq:total-jump-bound}
 \sum_\tau|\Delta J(\tau)|
 \le C(2n_0^3+n_0^2n_2+n_0^2n_1)
 \le C\mathcal S(A).
\end{equation}
The last step uses $a^2b\le(2a^3+b^3)/3$.
In particular, the estimate has no dependence on how many distinct
scales occur.

\paragraph{Telescoping the heat estimate.}
List the distinct squared endpoints as
$t_0=s_-<t_1<\cdots<t_m=s_+$.
On every open interval $(t_k,t_{k+1})$ the array is fixed, so
Lemma~\ref{lem:plane-dissipation} applies.  Its fixed-array continuity
and Gaussian domination also justify the one-sided endpoint limits.
Integrating on these intervals and summing gives
\begin{align*}
 c_*\int_{s_-}^{s_+}\int_{\Pi_0}\sum_vV_v\,\dd\pi\,\dd s
 &\le\sum_{k=0}^{m-1}\bigl[J(t_k+)-J(t_{k+1}-)\bigr]\\
 &=J(s_-+)-J(s_+-)
       +\sum_{k=1}^{m-1}\bigl[J(t_k+)-J(t_k-)\bigr]\\
 &\le J(s_-+)+\sum_{k=1}^{m-1}|\Delta J(t_k)|.
\end{align*}
Here the terminal energy was discarded by nonnegativity.  The initial
right-hand mask has entries bounded by those of the original array,
and all its variances are at least $h^2$.  Lemma~\ref{lem:initial-energy}
therefore gives $J(s_-+)\le C\mathcal S(A)$.
Together with \eqref{eq:total-jump-bound}, this proves
\eqref{eq:masked-dissipation}.  The initial opening is accounted for
by the right-hand initial energy; a final closing jump is unnecessary.
Simultaneous endpoints have already been combined in the list of
distinct scales, and their values at the single endpoint do not affect
the scale integral.
\end{proof}

Figure~\ref{fig:mask-energy} illustrates the masks and the energy accounting.
\begin{figure}[htbp]
\centering
\begin{tikzpicture}[x=1cm,y=1cm,font=\small,
  active/.style={line width=2pt,draw=blue!65!black},
  event/.style={densely dotted,draw=black!24},
  endpoint/.style={circle,draw=blue!65!black,fill=white,
    line width=.65pt,inner sep=0pt,minimum size=4pt},
  endpoint label/.style={font=\footnotesize,fill=white,inner sep=2pt},
  energy/.style={line width=1pt,draw=black},
  jump/.style={-{Stealth[length=3pt]},densely dashed,
    line width=.8pt,draw=red!65!black}]
  \foreach \t in {2.3,3.8,6.4,9.0,11.6}
    \draw[event] (\t,.1)--(\t,6.15);

  \node[anchor=west] at (1.9,6.65) {Output-dependent scale intervals};
  \foreach \yy in {5.8,4.8,3.8}
    \draw[black!22] (1.95,\yy)--(12.25,\yy);
  \node[anchor=east] at (1.7,5.8) {$(i_1,j_1)$};
  \node[anchor=east] at (1.7,4.8) {$(i_2,j_2)$};
  \node[anchor=east] at (1.7,3.8) {$(i_3,j_3)$};
  \draw[active] (2.3,5.8)--(6.4,5.8);
  \draw[active] (3.8,4.8)--(9.0,4.8);
  \draw[active] (6.4,3.8)--(11.6,3.8);
  \foreach \xx/\yy in {2.3/5.8,6.4/5.8,3.8/4.8,9.0/4.8,6.4/3.8,11.6/3.8}
    \node[endpoint] at (\xx,\yy) {};
  \node[endpoint label,above] at (2.3,5.91) {$\varepsilon_{i_1j_1}^{\,2}$};
  \node[endpoint label,above] at (6.4,5.91) {$\rho_{i_1j_1}^{\,2}$};
  \node[endpoint label,above] at (3.8,4.91) {$\varepsilon_{i_2j_2}^{\,2}$};
  \node[endpoint label,above] at (9.0,4.91) {$\rho_{i_2j_2}^{\,2}$};
  \node[endpoint label,above] at (6.4,3.91) {$\varepsilon_{i_3j_3}^{\,2}$};
  \node[endpoint label,above] at (11.6,3.91) {$\rho_{i_3j_3}^{\,2}$};
  \draw[-{Stealth[length=4pt]},black!65] (1.95,3.35)--(12.6,3.35)
    node[right] {$s$};

  \node[anchor=west,fill=white,inner sep=2pt] at (1.9,2.9) {Integrated energy};
  \draw[-{Stealth[length=4pt]},black!65] (1.95,.1)--(12.6,.1)
    node[right] {$s$};
  \node[anchor=east] at (1.7,1.7) {$J(s)$};
  \draw[energy] (2.3,2.2) .. controls (2.9,2.05) and (3.3,1.91) .. (3.8,1.82);
  \draw[jump] (3.8,1.82)--(3.8,2.38);
  \draw[energy] (3.8,2.38) .. controls (4.7,2.13) and (5.4,1.85) .. (6.4,1.64);
  \draw[jump] (6.4,1.64)--(6.4,1.97);
  \draw[energy] (6.4,1.97) .. controls (7.2,1.75) and (8.1,1.37) .. (9.0,1.18);
  \draw[jump] (9.0,1.18)--(9.0,.78);
  \draw[energy] (9.0,.78) .. controls (9.9,.61) and (10.7,.43) .. (11.6,.32);
  \node[below] at (2.3,.02) {$s_-$};
  \node[below] at (11.6,.02) {$s_+$};
\end{tikzpicture}
\caption{Three selected output pairs and a schematic integrated energy,
not an exact heat-flow solution. The entry $A_0^s(i,j)$ is retained only
on $(\varepsilon_{ij}^2,\rho_{ij}^2)$. Between successive event scales,
$J(s)$ is nonincreasing by the fixed-array dissipation inequality.
Dashed arrows connect its one-sided limits at events; the values at the
events themselves are not shown. Jumps may have either sign, and
endpoints may coincide, as illustrated
by $\rho_{i_1j_1}^2=\varepsilon_{i_3j_3}^2$. The bound
$\sum_\tau |\Delta J(\tau)|\le C\mathcal S(A)$, together with the
initial-energy bound, controls the accumulated dissipation independently
of the number of event scales.}
\label{fig:mask-energy}
\end{figure}

